\documentclass[11pt]{article}
\usepackage[a4paper,margin=27mm]{geometry}
\usepackage{amsmath,amssymb,amsthm,mathtools,hyperref}
\newtheorem{theorem}{Theorem}
\newtheorem{lemma}[theorem]{Lemma}
\newtheorem{corollary}[theorem]{Corollary}
\newtheorem{proposition}[theorem]{Proposition}
\newcommand{\At}{\operatorname{At}}
\newcommand{\cS}{\mathcal S}
\newcommand{\cP}{\mathcal P}
\title{No causality gap for atomic capacity\\on a locally finite infinite tree}
\author{Research note developed with Jan Mikul\'{\i}k}
\date{10 October 2026}
\begin{document}
\maketitle
\begin{abstract}
For the transcript laws of all finite-horizon deterministic policies on a
locally finite controlled tree, define $a_K$ to be the smallest mass carried
by the $K$ most likely transcripts, uniformly over policies and horizons.
We prove that $q=\lim_K a_K$ is exactly the maximum total atomic mass of an
exact causal representation. The maximum is attained and equals the
corresponding maximum without causality constraints. Thus the unavoidable
nonatomic mass is exactly $1-q$. The proof uses finite causal greedy
extraction, a quantitative concentration estimate, and a diagonal limit
with an explicitly realized residual flow. A further compactness argument
shows that bounded finite-horizon optimal entropies give one optimal
countable seed for the entire infinite tree, with no extra entropy cost.
\end{abstract}

\paragraph{Status and scope.}
This is a mathematical proof conditional on the finite causal extraction
and minimax theorem of the finite-Shannon-theorem project. It is not a Lean
certificate, and publication priority is not claimed. Infinite families of
ordinary couplings and their uniform concentration condition are already
treated by Li~\cite{Li}. The assertion examined here is attainment of the
\emph{same atomic capacity by compatible strategies}, including $q<1$.
This does not determine the finite Shannon entropy gap $G_*$.

\section{Objects and statement}
Let $T$ be a rooted controlled history tree. At every finite history $h$
there are finitely many available actions, each with a finite nonempty
output alphabet and a prescribed probability kernel. We assume the
alphabets are finite at every node, not uniformly bounded. Terminal nodes
can be padded by a forced action and output. Zero-probability outputs may
be deleted; deterministic singleton alphabets are permitted.

Write $p(h)$ for the product of output probabilities along $h$.
For every available action $a$,
\[
 p(\varnothing)=1,\qquad \sum_y p(hay)=p(h).
\]
At each finite depth there are finitely many histories. A deterministic
response strategy selects one output at every history--action pair.
The space $\cS$ of all such strategies is a countable product of finite
discrete spaces, hence compact and metrizable. For a history $h$, let
$E_h\subseteq\cS$ be the set of strategies agreeing with its output
choices. This is a clopen cylinder, and
\[
 E_h=\bigsqcup_y E_{hay}\quad\text{for every action }a.
\]
A Borel probability measure $\mu$ on $\cS$ is an exact causal
representation if $\mu(E_h)=p(h)$ for every finite history $h$.
This is equivalent to exactness for every deterministic adaptive policy
at every finite horizon, and therefore for the corresponding infinite
path laws.

Let $P_{n,\pi}$ be the finite transcript law at horizon $n$ under policy
$\pi$. Put
\begin{align}
 A_K(P)&=\max_{D:\,|D|\le K}P(D),\\
 a_K(T)&=\inf_{n,\pi} A_K(P_{n,\pi}),\qquad
 q(T)=\lim_{K\to\infty}a_K(T).\label{eq:q}
\end{align}
The limit exists by monotonicity and belongs to $[0,1]$.
For a Borel probability on a compact metrizable space define its total
atomic mass by
\[
 \At(\mu)=\sum_{s:\,\mu(\{s\})>0}\mu(\{s\}).
\]
There are at most countably many terms.

Ordinary couplings below are Borel probability laws on the countable
product of the transcript spaces indexed by all pairs $(n,\pi)$, with
the specified coordinate marginals. They need not be compatible with a
single strategy or even consistent between horizons.

\begin{theorem}[Atomic capacity]\label{th:main}
For every such tree,
\[
 \boxed{\quad
 \max_{\mu\ \mathrm{exact\ causal}}\At(\mu)
 =\max_{\Gamma\ \mathrm{ordinary}}\At(\Gamma)
 =q(T).\quad}
\]
Both maxima are attained. In particular, an exact countably supported
causal seed for the entire infinite tree exists if and only if $q(T)=1$.
If $q(T)<1$, every exact causal representation has nonatomic mass at
least $1-q(T)$, and a representation attaining this minimum exists.
\end{theorem}

\section{The finite concentration estimate}
We use the following part of the finite theorem~\cite{repo}. After a
finite number of causal greedy extractions, the residual cylinder masses
form a nonnegative controlled flow. If its root mass is $R$, its next
extractable strategy weight $w$ equals the smallest, over policies, of
the largest residual transcript mass. Extraction terminates on each
finite tree and gives an exact finite seed.

\begin{lemma}\label{lem:concentration}
Suppose every policy on a finite controlled tree has $K$ transcripts
carrying total mass at least $a$. The first $M$ greedy atoms, padded with
zeros if extraction has ended, have total mass at least
\[
 a\bigl[1-(1-1/K)^M\bigr].\label{eq:head}
\]
Consequently the same is true of the $M$ largest seed atoms.
\end{lemma}
\begin{proof}
Let $R_m$ be the residual root mass after $m$ steps, with $R_0=1$.
The minimax identity supplies a policy all of whose residual transcript
masses are at most the next weight $w_{m+1}$. On this policy select at
most $K$ original transcripts of mass at least $a$. The residual mass
outside this set is at most $1-a$, so
\[
 R_m-(1-a)\le K w_{m+1}.
\]
Since $R_{m+1}=R_m-w_{m+1}$,
\[
 R_{m+1}-(1-a)\le(1-1/K)(R_m-(1-a)).
\]
Iteration gives $R_M\le1-a+a(1-1/K)^M$. The same bound holds after
termination. Subtract from one. For $K=1$ and $M\ge1$ the displayed
formula has its usual value $a$.
\end{proof}

\section{Proof of atomic capacity}
\begin{proof}[Proof of Theorem~\ref{th:main}]
\emph{Upper bound, even without causality.}
Take any $m$ atoms of an ordinary coupling, of total mass $b$. Their
images in any one transcript coordinate lie in at most $m$ transcripts,
so $b\le A_m(P_{n,\pi})$ for every $(n,\pi)$. Hence $b\le a_m\le q$.
Exhausting all its atoms gives $\At(\Gamma)\le q$. The same argument
applies directly to a causal representation by projecting each atom
through a policy.

\emph{Limits of weights and strategies.}
For each horizon $n$, take a finite greedy causal seed, order its weights
as $w^{(n)}_1\ge w^{(n)}_2\ge\cdots\ge0$, and pad by zeros.
Extend each finite strategy arbitrarily to an element $s^{(n)}_i$ of
$\cS$. By compactness and diagonal selection there is one subsequence
$n_j\to\infty$ such that, for every $i$,
\[
 w_i^{(n_j)}\to w_i,\qquad s_i^{(n_j)}\to s_i.
\]
Set $\nu=\sum_iw_i\delta_{s_i}$ and $A=\sum_iw_i\le1$; repeated
limit strategies are simply merged as measures. For fixed $K,M$,
Lemma~\ref{lem:concentration}, with $a=a_K(T)$, gives
\[
 \sum_{i=1}^M w_i
 \ge a_K(T)\bigl[1-(1-1/K)^M\bigr].
\]
First let $M\to\infty$ and then $K\to\infty$. It follows that
$A\ge q$.

\emph{No invalid passage of infinitely many atoms through a limit.}
For fixed $h$ and $M$, convergence of strategies makes membership in
the clopen set $E_h$ eventually constant for each of the first $M$
indices. For large $j$, the finite seed is exact at $h$. Thus
\[
 \sum_{i=1}^M w_i\mathbf1_{E_h}(s_i)
 =\lim_j\sum_{i=1}^M w_i^{(n_j)}\mathbf1_{E_h}(s_i^{(n_j)})
 \le p(h).
\]
Monotone convergence in $M$ proves $\nu(E_h)\le p(h)$. Define
\[
 r(h)=p(h)-\nu(E_h)\ge0.
\]
Since the $E_{hay}$ partition $E_h$ for each action,
$\sum_y r(hay)=r(h)$ and $r(\varnothing)=1-A$.

\emph{Realization of the missing mass.}
If $A<1$, assign, independently at every history--action pair, an output
with probabilities $r(hay)/r(h)$ whenever $r(h)>0$, and any distribution
when $r(h)=0$. The countable product law $\lambda$ of these response
choices satisfies
\[
 \lambda(E_h)=r(h)/(1-A).
\]
Indeed this follows by induction along $h$ from the displayed ratios;
after a zero residual prefix both sides vanish. Hence
$\mu=\nu+(1-A)\lambda$ is exact. If $A=1$, simply set $\mu=\nu$.
As $\nu$ is atomic of mass $A$, $\At(\mu)\ge A\ge q$.
The upper bound already proved gives $\At(\mu)=A=q$. In particular,
the residual law in this construction has no atoms when $q<1$.

Finally map $s\in\cS$ to its full tuple of finite-policy transcripts.
This is a measurable map into the countable product of finite spaces.
The image of $\mu$ is an ordinary coupling and has atomic mass at least
$q$, since images of atoms remain atoms. The upper bound makes its
atomic mass exactly $q$. This proves attainment on both sides.
\end{proof}

\section{Information spectrum and entropy}
All logarithms in this section have base two. Define
\[
 F_n(t)=\min_\pi P_{n,\pi}\{-\log P_{n,\pi}(X)\le t\},\quad
 F_\infty(t)=\inf_n F_n(t),\quad t\ge0.
\]
Refining a transcript cannot reduce its information, and policies extend
to longer horizons. Therefore $F_n(t)$ decreases with $n$.
For an individual finite law $P$, writing its information CDF as $F_P$,
\[
 F_P(t)\le A_{\lfloor2^t\rfloor}(P),\qquad
 A_K(P)\le F_P(t)+K2^{-t}.
\]
Taking infima over all finite-policy laws proves
\[
 \boxed{q(T)=\lim_{t\to\infty}F_\infty(t).}
\]
Thus atomic capacity is precisely the mass that does not escape to
infinite information in the uniform policy envelope. This statement
does not require treating a possibly defective $F_\infty$ as a
probability CDF.

Let $C_n$ be the minimum entropy of an exact seed for the depth-$n$
tree. Let $C_\infty$ be the infimum entropy among countably supported
exact seeds for the full tree, with value $+\infty$ if there is no such
seed. Set
\[
 L_n=\int_0^\infty(1-F_n(t))\,dt,\quad
 L_\infty=\sup_nL_n,\quad c_*=(1+\log e)/2.
\]

\begin{theorem}[No extra entropy at the projective limit]
In the extended real numbers,
\[
 \boxed{C_\infty=\sup_n C_n.}
\]
If the right-hand side is finite, the infimum on the left is attained.
Consequently the finite Shannon theorem implies
\[
 L_\infty\le C_\infty\le L_\infty+c_*.
\]
In particular, $L_\infty<\infty$ is equivalent to the existence of one
finite-entropy exact causal seed for the entire infinite tree.
\end{theorem}
\begin{proof}
Restriction of a seed proves $C_n\le C_{n+1}\le C_\infty$. If
$B=\sup_n C_n=+\infty$, this already proves the equality.
Assume $B<\infty$ and choose a minimizing finite seed at each horizon.
Such minimizers exist: one can minimize entropy over the compact
finite-dimensional polytope of laws on finite strategies, merging
identical strategies if necessary. Order and pad weights as above.
For $i>K$, $w_i^{(n)}\le1/(K+1)$, whence
\[
 \sum_{i>K}w_i^{(n)}\le \frac{B}{\log(K+1)}.\tag{*}
\]
Diagonal compactness gives limiting weights and strategies. The uniform
tail bound (*) forces $\sum_iw_i=1$. The cylinder argument in the
previous proof gives $\nu(E_h)\le p(h)$. Since both are probability
flows, equality follows at each depth by summing over the leaves of
any policy following $h$, or by passing the uniformly small tails
directly through the limit.

The countable label seed $i$ with masses $w_i$ realizes these strategies.
With $0\log0=0$, continuity of each summand and Fatou's lemma give
\[
 H(w)=\sum_i-w_i\log w_i
 \le\liminf_j H(w^{(n_j)})=B.
\]
Together with restriction this yields $C_\infty=B$, with attainment.
Finally the finite theorem gives $L_n\le C_n\le L_n+c_*$;
take suprema. No continuity of Shannon entropy on an unrestricted
countable simplex is being assumed.
\end{proof}

\paragraph{Three regimes.}
If $q<1$, every exact representation has an unavoidable nonatomic
component. If $q=1$ but $L_\infty=\infty$, a countable exact seed exists
but all such seeds have infinite entropy. If $L_\infty<\infty$, an
optimal finite-entropy countable seed exists. Countable support is not
the same as finite support.

\section{A strict quantifier separation}
Consider a deterministic waiting chain. At waiting stage $m\ge1$ the
controller can either wait for the next stage or trigger an experiment.
Triggering at stage $m$ produces $m$ independent fair bits and then
stops; there are no choices during the experiment. A policy that never
triggers has a deterministic infinite transcript. The tree is locally
finite, with at most two actions and at most two outputs at every node.

Every individual deterministic policy therefore has a finite-support
full-path law: either one path, or $2^m$ paths for some finite $m$.
Nevertheless, for each fixed $K$, choosing arbitrarily large $m$ yields
\[
 a_K\le K2^{-m},\qquad a_K=0,\qquad q=0.
\]
Hence \emph{every exact common representation is atomless}, even
without a causality requirement. An elementary direct proof is that
an atom of mass $w>0$ would have to map into a transcript of mass
$2^{-m}<w$ for a sufficiently large triggered experiment.
This example illustrates why separate finite-support simulations do
not imply one discrete seed. It is not claimed as a new counterexample
in the general theory of ordinary coupling.

For any prescribed $q_0\in(0,1)$, first choose a deterministic absorbing
branch with probability $q_0$ and the waiting construction with
probability $1-q_0$. Then $a_K=q_0$ for every $K\ge1$: the absorbing
transcript gives the lower bound, and triggering an arbitrarily long
experiment gives the upper bound. Thus the full interval of atomic
capacities is realized even with binary branching.

\paragraph{What the result does not say.}
The diagonal argument is an existence proof, not a uniform algorithm
for an arbitrary computable infinite tree. It does not bound the time
needed to sample the seed, determine $G_*$, prove a new asymptotic
coefficient, or assert that this infinite extension has been formally
verified. Its structural conclusion is the exact equality of ordinary
and causal atomic capacities, including attainment and partial atomic
mass, and the absence of a further entropy charge when finite-horizon
optima stay bounded.

\begin{thebibliography}{9}
\bibitem{repo} J. Mikul\'{\i}k, \emph{An exact causal seed on a finite
controlled tree}, finite-Shannon-theorem repository. The finite
minimax/extraction and Shannon statements are the inputs used here.
\url{https://github.com/jendamikulik/finite-Shannon-theorem}.
\bibitem{Li} C. T. Li, \emph{Efficient Approximate Minimum Entropy
Coupling of Multiple Probability Distributions}, IEEE Transactions on
Information Theory 67(8), 5259--5268, 2021.
\url{https://arxiv.org/abs/2006.07955}.
Definition 6 and Corollary 9 treat uniform concentration and ordinary
discrete couplings for infinite families; this prior work must be
distinguished from causal compatibility on a history tree.
\end{thebibliography}
\end{document}
